\documentclass[11pt,a4paper]{article}

\usepackage[T1]{fontenc}
\usepackage[utf8]{inputenc}
\usepackage[spanish,es-noquoting]{babel}
\usepackage{mathptmx}\usepackage[scaled=.9]{helvet}
\usepackage{amsmath,amssymb}
\usepackage{graphicx}
\usepackage{booktabs}
\usepackage{longtable}
\usepackage{xcolor}
\usepackage[margin=2.3cm]{geometry}
\usepackage{tikz}
\usetikzlibrary{arrows.meta,positioning,fit,calc}
\usepackage{float}
\usepackage{enumitem}
\usepackage[colorlinks=true,linkcolor=azulq,urlcolor=azulq]{hyperref}
\usepackage{caption}

\definecolor{azulq}{HTML}{1F3B59}
\definecolor{cuantico}{HTML}{5347C4}
\definecolor{clasico}{HTML}{A8651F}
\definecolor{fiel}{HTML}{2F7A5C}
\definecolor{gris}{HTML}{6B6B76}

\captionsetup{font=small,labelfont={bf,color=azulq}}
\newcommand{\M}[1]{\texttt{#1}}

\title{\bfseries Guía del código\\[.3em]
  \large Qué hace cada archivo, por dónde pasa una corrida,\\
  y dónde buscar cuando algo no cuadra}
\author{Andrés García Rivera}
\date{\today}

\begin{document}
\maketitle

\begin{abstract}
\noindent
El repositorio son unas $15\,000$ líneas repartidas en doce archivos, y no
hace falta entenderlas todas para trabajar con él. Esta guía separa lo que
\emph{hay que} entender de lo que solo hay que saber \emph{dónde está}.

\medskip\noindent
El inventario de esta guía se regenera con
\M{python docs/guia\_codigo/inventario.py}, leyendo el árbol de sintaxis de
cada módulo. Si los números de aquí no coinciden con los de ese comando, los
buenos son los del comando.
\end{abstract}

\tableofcontents
\bigskip

% =====================================================================
\section{El mapa}
% =====================================================================

Doce archivos, dos grupos con propósitos distintos.

\subsection{El motor — lo que corre la ciencia}

\begin{center}
\small
\begin{tabular}{@{}lrp{7.4cm}@{}}
\toprule
\textbf{Módulo} & \textbf{líneas} & \textbf{Qué hace / cuándo lo abres} \\
\midrule
\M{cosmo\_core.py} & 1\,757
  & La \textbf{física y los datos}. Modelos cosmológicos, carga de
    catálogos, $\chi^2$, priors, $\hat{R}$, ESS. \emph{Lo abres si cambias
    un modelo, un dataset o una estadística.} \\
\addlinespace
\M{cosmo\_noise.py} & 864
  & El \textbf{eje de ruido}. Define los peldaños y construye el simulador
    de cada uno. \emph{Lo abres si tocas el ruido.} \\
\addlinespace
\M{cosmo\_modular\_quantum.py} & 4\,347
  & Los \textbf{muestreadores}: QMCMC y QVMC, con sus componentes
    conmutables. Es el archivo más grande y el corazón del proyecto. \\
\addlinespace
\M{cosmo\_genetic\_optimizers.py} & 3\,369
  & Los \textbf{genéticos}: CGA y QGA con sus tres operadores
    conmutables. \\
\addlinespace
\M{cosmo\_hpc\_runner.py} & 2\,209
  & El \textbf{orquestador}. Reparte las tareas de una campaña entre
    núcleos, arma las carpetas y junta los resultados. No tiene física
    dentro. \\
\addlinespace
\M{qpu\_cosmo\_samplers.py} & 1\,465
  & La ruta de \textbf{hardware real} (IBM Quantum). Se mantiene intacta a
    propósito; el simulador no la toca. \\
\addlinespace
\M{cosmo\_profiling.py} & 374
  & Mide memoria y tiempo de pared. Lo usa el runner, no tú. \\
\addlinespace
\M{data\_manifest.py} & 158
  & Sumas de verificación de los datos, para reproducibilidad. \\
\bottomrule
\end{tabular}
\end{center}

\subsection{El análisis — lo que lee resultados ya obtenidos}

Ninguno de estos corre nada. Los cuatro leen los CSV de una campaña
terminada. Esto es deliberado: analizar no debe costar otra campaña.

\begin{center}
\small
\begin{tabular}{@{}lrp{7.4cm}@{}}
\toprule
\textbf{Script} & \textbf{líneas} & \textbf{Qué pregunta contesta} \\
\midrule
\M{triage\_campana.py} & 461
  & «¿Esta campaña está sana?» Revisa tareas muertas, ajuste, procedencia y
    las celdas fieles. \emph{Es lo primero que corres al bajar
    resultados.} \\
\addlinespace
\M{comparar\_semillas.py} & 403
  & «¿Esta diferencia sobrevive a cambiar la semilla?» \\
\addlinespace
\M{graficas\_ruido.py} & 510
  & «¿Cuánto mueve el ruido al resultado, y crece con el tamaño del
    registro?» \\
\addlinespace
\M{comparar\_algoritmos.py} & 398
  & «¿Cuál método es mejor, y con qué número se mide eso?» \\
\bottomrule
\end{tabular}
\end{center}

% =====================================================================
\section{Por dónde pasa una corrida}
% =====================================================================

Éste es el recorrido completo, de la línea de comandos al CSV. Si entiendes
este diagrama, entiendes la arquitectura.

\begin{figure}[H]
\centering
\begin{tikzpicture}[
  font=\small,
  n/.style={draw,rounded corners=2pt,align=center,inner sep=5pt,
            minimum height=8mm,text width=34mm,line width=.5pt},
  orq/.style={n,draw=gris,fill=gris!8},
  fis/.style={n,draw=clasico,fill=clasico!8},
  alg/.style={n,draw=cuantico,fill=cuantico!10},
  sal/.style={n,draw=fiel,fill=fiel!10},
  fl/.style={-{Stealth[length=4pt]},draw=gris,line width=.6pt},
]
\node[orq] (cli) {\M{cosmo\_hpc\_runner.py}\\[-2pt]{\scriptsize la campaña}};
\node[orq,below=5mm of cli] (tarea)
  {una \textbf{Task}\\[-2pt]{\scriptsize modelo × rejilla × ruido}};
\node[fis,below=5mm of tarea] (noise)
  {\M{cosmo\_noise}\\[-2pt]{\scriptsize fija el peldaño}};
\node[fis,below=5mm of noise] (post)
  {\M{cosmo\_core.Posterior}\\[-2pt]{\scriptsize modelo + datos + prior}};
\node[alg,below=5mm of post] (met)
  {el método\\[-2pt]{\scriptsize QMCMC / QVMC / QGA}};
\node[sal,below=5mm of met] (csv)
  {\M{resultados\_config.csv}\\[-2pt]{\scriptsize una fila por peldaño}};
\node[sal,below=5mm of csv] (ana)
  {los cuatro scripts\\[-2pt]{\scriptsize de análisis}};
\foreach \i/\j in {cli/tarea,tarea/noise,noise/post,post/met,met/csv,csv/ana}
  \draw[fl] (\i) -- (\j);

\node[text width=6.3cm,right=8mm of tarea,align=left,font=\scriptsize]
  {Una \textbf{tarea} es una carpeta:
   \M{samplers\_lcdm\_nqpp3\_noise-readout}. El nombre lleva la familia, el
   modelo, la rejilla y el peldaño — y por eso el análisis puede
   reconstruir la procedencia aunque una columna del CSV esté mal.};
\node[text width=6.3cm,right=8mm of post,align=left,font=\scriptsize]
  {\textbf{Aquí y solo aquí} entran los 1\,099 datos. \M{Posterior.log\_prob}
   es la única puerta por la que el catálogo toca el cálculo.};
\node[text width=6.3cm,right=8mm of met,align=left,font=\scriptsize]
  {El método corre \emph{toda su escalera} en una tarea: el clásico y los
   peldaños cuánticos, uno tras otro, compartiendo semilla. Por eso una
   celda fiel se puede comparar dígito a dígito.};
\end{tikzpicture}
\caption{De la línea de comandos al CSV. El color marca el papel: gris
orquestación, ocre física, índigo el algoritmo conmutable, verde la salida.}
\end{figure}

\paragraph{Lo que hay que retener.} El runner no sabe de física; la física
no sabe de cuántica; el método no sabe de datos (los pide por
\M{log\_prob}). Esa separación es lo que permite cambiar un componente sin
tocar el resto — y es la razón de que el marco de ablación sea posible.

% =====================================================================
\section{Los cinco objetos que sí hay que entender}
% =====================================================================

Todo lo demás es andamiaje. Estos cinco son el proyecto.

\subsection{\M{Posterior} — la función objetivo}
\emph{en \M{cosmo\_core.py}}

Junta modelo $+$ dataset $+$ prior y expone lo único que los métodos
necesitan:

\begin{center}
\small
\begin{tabular}{@{}ll@{}}
\toprule
\M{post.log\_prob(theta)} & log-posterior en un punto; \textbf{aquí entran los datos} \\
\M{post.log\_prob\_batch(thetas)} & lo mismo vectorizado — es lo que usa el MCMC \\
\M{post.chi2(theta)} & $(\chi^2, n_{\text{datos}})$ \\
\M{post.model.n\_params} & $d$, y de ahí sale el ancho de los circuitos \\
\M{post.model.sample\_box} & la caja de búsqueda por parámetro \\
\bottomrule
\end{tabular}
\end{center}

\noindent
Si algo sale raro en la física, se depura aquí, sin cuántica de por medio.

\subsection{\M{NoiseSpec} — el peldaño de ruido}
\emph{en \M{cosmo\_noise.py}}

\M{NoiseSpec.from\_level('readout')} devuelve el objeto que sabe construir
el simulador correspondiente. Dos detalles que ya mordieron:

\begin{itemize}[leftmargin=1.6em]
  \item \M{simulator\_kwargs(counts\_route=...)} decide \textbf{quién}
    aplica el canal de lectura. Equivocarlo produce doble conteo o
    desaparición silenciosa del error, así que la decisión no se deja al
    que llama.
  \item Cada módulo tiene su \textbf{propio} \M{NOISE} global y su propio
    \M{set\_noise}. Olvidar uno fue el bug \M{[B-PROV-GEN]}: el genético
    corría con ruido y lo archivaba como ideal.
\end{itemize}

\subsection{\M{QMCMCModular} — la cadena}
\emph{en \M{cosmo\_modular\_quantum.py}}

Metropolis--Hastings escrito a mano, \textbf{no} \M{emcee}. Está así a
propósito: para intercambiar componentes hay que ser dueño de cada línea
del núcleo de transición, y \M{emcee} es una caja negra que además haría
imposible la comparación pareada.

Dos componentes conmutables: \M{proposal} y \M{acceptance}. El bucle está
vectorizado sobre cadenas — un solo \M{log\_prob\_batch} por paso, porque
la verosimilitud es el costo dominante.

\subsection{\M{QVMCModular} — el variacional}
\emph{en \M{cosmo\_modular\_quantum.py}}

Representa la posterior discretizada en $2^{d\cdot n_{qpp}}$ puntos como
$|\psi(\varphi)|^2$ y minimiza $\mathrm{KL}(Q_\varphi\|P)$. Tres
componentes: \M{sampling}, \M{training}, \M{normalization}.

Los métodos que importan: \M{build\_target()} arma $P$ sobre la rejilla;
\M{\_kl\_batch()} evalúa el ansatz y devuelve la KL; \M{train()} optimiza;
\M{sample()} extrae el estimador.

\subsection{\M{CGA} y \M{QGA} — los genéticos}
\emph{en \M{cosmo\_genetic\_optimizers.py}}

Comparten \M{GeneticEvolver}, así que el clásico y el cuántico recorren el
mismo bucle y solo cambian los operadores: \M{q\_init}, \M{q\_mutation},
\M{q\_crossover}. Los tres circuitos se transpilan \textbf{una sola vez} en
\M{\_build\_circuits()} y después solo se ligan ángulos.

% =====================================================================
\section{Dónde vive cada cosa}
% =====================================================================

La tabla que se consulta cuando se sabe qué cambiar pero no dónde.

\begin{center}
\small
\begin{longtable}{@{}p{7.1cm}p{7.6cm}@{}}
\toprule
\textbf{Quiero\ldots} & \textbf{Está en\ldots} \\
\midrule
\endhead
añadir un modelo cosmológico
  & \M{cosmo\_core.py}, el diccionario \M{MODELS} \\
añadir o cambiar un dataset
  & \M{cosmo\_core.py}, los \M{chi2\_*} y el cargador \\
cambiar un prior
  & \M{cosmo\_core.Posterior.log\_prior} \\
cambiar el umbral de convergencia
  & \M{cosmo\_core.RHAT\_THRESHOLD} (hoy $1.01$) \\
\addlinespace
añadir un peldaño de ruido
  & \M{cosmo\_noise.NoiseSpec.from\_level} y \M{NAMED\_LEVELS} \\
cambiar las tasas de error
  & \M{cosmo\_noise.DEFAULT\_READOUT\_P}, \M{\_GATE\_P1}, \M{\_GATE\_P2} \\
\addlinespace
cambiar el circuito de propuesta
  & \M{cosmo\_modular\_quantum.build\_proposal\_circuit} \\
cambiar cómo se lee la propuesta
  & \M{QuantumProposalEngine.\_raw\_block\_amplitude} / \M{\_counts} \\
cambiar el ansatz variacional
  & \M{QVMCModular.\_build\_ansatz} \\
cambiar los operadores genéticos
  & \M{QGA.\_build\_circuits} (init, mutación, cruce) \\
\addlinespace
cambiar qué columnas salen al CSV
  & \M{cosmo\_modular\_quantum.csv\_fields\_for\_model} y
    \M{csv\_provenance} \\
cambiar cómo se reparten las tareas
  & \M{cosmo\_hpc\_runner.py}, la clase \M{Task} y el planificador \\
\bottomrule
\end{longtable}
\end{center}

% =====================================================================
\section{Los marcadores del código}
% =====================================================================

Al leer los módulos aparecen etiquetas entre corchetes. No son ruido: cada
una marca una decisión que costó encontrar, y casi todas nacieron de un bug
medido.

\begin{center}
\small
\begin{tabular}{@{}llp{8.1cm}@{}}
\toprule
\textbf{Etiqueta} & \textbf{veces} & \textbf{Qué marca} \\
\midrule
\M{[NOISE]}     & 28 & consecuencia del eje de ruido en ese punto \\
\M{[FIX]}       & 24 & corrección de un error anterior \\
\M{[OPT]}       & 20 & optimización de rendimiento y por qué era necesaria \\
\M{[B-\ldots]}  & --- & un bug concreto, con su medición \\
\bottomrule
\end{tabular}
\end{center}

\noindent
Los \M{[B-\ldots]} son los que más vale la pena leer, porque cada uno
documenta una trampa real. Los que más aparecen:

\begin{center}
\small
\begin{tabular}{@{}lp{10.4cm}@{}}
\toprule
\M{[B-REFINE]} & el $\chi^2$ refinado del genético no distingue peldaños;
                 \M{chi2\_grid} sí \\
\M{[B-PROV]}   & el CSV debe decir con qué ruido y semilla se obtuvo cada
                 fila \\
\M{[B-BUDGET]} & presupuestos igualados entre el clásico y el cuántico \\
\M{[B-ESSCOMP]}& el ESS del variacional es el número de disparos, no
                 calidad \\
\M{[B-EMPTY]}  & una rejilla sin soporte del prior daba una KL
                 engañosamente buena \\
\M{[B-PROV-GEN]}& el genético archivaba toda corrida ruidosa como ideal \\
\M{[B-RZZ]}    & el peldaño \M{full} mataba el transpilador del cruce \\
\M{[B-EPS15]}  & un \M{+1e-15} rompía la igualdad exacta de una celda
                 fiel \\
\bottomrule
\end{tabular}
\end{center}

% =====================================================================
\section{Cómo leerlo sin ahogarse}
% =====================================================================

Cuatro mil líneas seguidas no se leen. Orden sugerido, de menos a más:

\begin{enumerate}[leftmargin=1.8em]
  \item \textbf{Corre la traza.}
    \M{python docs/circuitos/traza\_un\_paso.py}. Ver los números cruzar
    las fronteras vale más que leer el bucle.
  \item \textbf{Lee \M{cosmo\_core.Posterior}} completo. Son pocas líneas
    y es la única puerta de los datos.
  \item \textbf{Lee un método entero, no los tres.} Empieza por
    \M{QMCMCModular}: es el más corto y el que tiene la celda fiel con
    demostración.
  \item \textbf{Lee los \M{[B-\ldots]}} del módulo que te toque. Son los
    comentarios largos, y cada uno te ahorra repetir un error.
  \item \textbf{Deja el runner para el final}, y solo si vas a cambiar
    cómo se lanza una campaña. No tiene física.
  \item \textbf{No leas \M{qpu\_cosmo\_samplers.py}} salvo que vayas a
    correr en hardware. Es una ruta paralela y se mantiene aparte a
    propósito.
\end{enumerate}

\paragraph{Y una regla práctica.} Si algo del resultado no cuadra, el orden
de sospecha es: primero los datos y el prior (\M{cosmo\_core}), después la
procedencia de la fila (¿qué peldaño y qué semilla dice el CSV?), después
el componente conmutado, y solo al final el circuito. En este proyecto la
mayoría de los bugs encontrados estuvieron en los primeros dos escalones,
no en la cuántica.

\end{document}
